\begin{claim}{C3.1}{Proved}{The real Bott--Chern Chern-class map and a global $G$}
Let $G$ be a $\bar\partial$-closed smooth $(0,1)$ representative of
a nonzero $g\in H^1(X,\OO_X)$.  Set
\[
\mathsf A=[\partial G]_{BC},\qquad
\mathsf B=[\bar\partial\bar G]_{BC}.
\]
These classes form a complex basis of $H^{1,1}_{BC}(X)$.
With the exponential $z\mapsto e^{2\pi iz}$ and the Chern convention
$c_{1,BC}=[iF/(2\pi)]$, the map of real Lie groups is
\begin{equation}\label{eq:C-Pic-map}
\Pic(X)=\{\exp(zg):z\in\C\}\longrightarrow H^{1,1}_{BC}(X),
\qquad \exp(zg)\longmapsto z\mathsf A+\bar z\mathsf B.
\end{equation}
It is injective and its image is the real subspace fixed by conjugation.

There is an operator-explicit global representative.  Normalize the
Fubini--Study form by $\int_{\Pone}\omega_{FS}=1$, put
$\theta=f^*\omega_{FS}$, choose any smooth Riemannian metric on $X$,
and let $\mathcal G_d$ be its Green operator for the real de Rham
Laplacian.  Then
\begin{equation}\label{eq:C-G-Green}
\alpha=d^*\mathcal G_d\theta,\qquad G_0=\alpha^{0,1}
\end{equation}
satisfy $d\alpha=\theta$, $\bar\partial G_0=0$,
$[G_0]\ne0$, and $\exp([G_0])=f^*\OO_{\Pone}(1)$.
\end{claim}
\begin{proof}
C1.1 shows that $[\partial G]_{\bar\partial}=d_1g\ne0$.
The conjugate class $\mathsf B$ is nonzero, and its Dolbeault image
is zero.  Thus $\mathsf A,\mathsf B$ are independent.  C2.1 gives
their completeness as a basis.

The exponential sequence and $H^1(X;\Z)=H^2(X;\Z)=0$ identify
$\Pic(X)$ with $H^1(X,\OO_X)=\C g$, with no period lattice.
On the smooth trivial line with unit frame $s$, define
$\bar\partial_zs=-2\pi izG\otimes s$.  This is integrable.
If $G=\bar\partial\phi_i$ locally, the holomorphic frames
$e_i=e^{2\pi iz\phi_i}s$ have transition functions
$e_j=e^{2\pi iz(\phi_j-\phi_i)}e_i$, which represent
$\exp(zg)$.  The metric-compatible Chern connection is
\[
\nabla_z=d-2\pi i(zG+\bar z\bar G),\qquad
F_z=-2\pi i(z\partial G+\bar z\bar\partial\bar G).
\]
This proves \eqref{eq:C-Pic-map}, including its sign.  The basis
proves injectivity and the assertion about its real image.

The form $\theta$ is closed and exact because $H^2(X;\R)=0$.
Hodge decomposition, together with $d\theta=0$ and the commutation
of $d$ with the Green operator, gives $dd^*\mathcal G_d\theta=
\theta$.  The operator preserves real forms.  Since $d\alpha$ has
type $(1,1)$, its $(0,2)$ component says
$\bar\partial G_0=0$.  If $G_0=\bar\partial u$, then
$\theta=\partial\bar\partial(u-\bar u)$.  A real multiple of
$\operatorname{Im}u$ would be a global plurisubharmonic function,
because $\theta$ is semipositive.  The maximum principle on compact
$X$ makes this function constant.  This contradicts the fact that
$\theta$ is nonzero over the regular locus of $f$.

Finally \eqref{eq:C-Pic-map} gives
$c_{1,BC}(\exp([G_0]))=[\theta]=c_{1,BC}(f^*\OO(1))$.
Injectivity gives the asserted line-bundle equality.  Formula
\eqref{eq:C-G-Green} is a global formula involving a specified
elliptic inverse; it is not a finite coordinate formula on the toric
and cyclic charts, and it depends on the auxiliary metric.
\end{proof}
